\documentclass[14pt]{article}
\usepackage{graphicx} 
\usepackage{moresize}
\usepackage{xcolor}
\pagecolor{black}
\usepackage{siunitx}
\usepackage{amsmath}
\usepackage{amssymb}
\color{white}
\usepackage{tikz}
\usepackage[letterpaper, margin=0.75in]{geometry}
\usepackage{background}
\usepackage{fancyhdr}
\usepackage{lastpage}

\pagestyle{fancy}
\fancyhf{}
\fancyfoot[C]{{\Large\textbf{\thepage}}/\pageref{LastPage}}
\renewcommand{\headrulewidth}{0pt}

\definecolor{goldColor}{HTML}{D4AF37}

\backgroundsetup{
    scale=1,
    color=goldColor,
    opacity=1,
    angle=0,
    position=current page.center,
    contents={
        \definecolor{goldColor}{HTML}{D4AF37}
        \tikz\draw[goldColor, line width=2pt, rounded corners=8pt] 
            (0,0) rectangle (\dimexpr\paperwidth-1in\relax, \dimexpr\paperheight-1in\relax);
    }
}


\begin{document}

\begin{center}
	\vspace*{30pt}
    {\Huge \textbf{ Global}} \\[15pt]
    {\HUGE \textbf{Quantum Mechanics}}\\[15pt]
    {\Huge \textbf{Challenge}}\\[50pt]
    {\Large \textcolor{goldColor}{\textbf{Qualification Round 2026}}}\\[100pt]
    {\huge \textcolor[HTML]{4DB8D7}{\textbf{Md Mokhdum Amin Fahim}}}\\[30pt]
    {\huge \textcolor{cyan}{Police Line Secondary School, Jashore}}\\[30pt]
    {\huge \textcolor[HTML]{E6E6FA}{\textbf{Bangladesh}}}\\[100pt]
    {\Large Date : \today}\\[15pt]
    {\Large Page Count : 03}
\end{center}

\newpage
\begin{center}
	\section*{\textcolor{goldColor}{\textbf{Problem A}}}
\end{center}
\textcolor{cyan}{
(a) Radio waves
; (b) Infrared
; (c) Ultraviolet
; (d) X-rays
; (e) Gamma rays
}\\\\
\section*{\centering \textcolor{goldColor}{\textbf{Problem B}}}

Given, 
\begin{itemize}
    \item Speed of light $c=3\times10^8$ m/s
    \item Wavelength of green light $\lambda_g=532$ nm $=5.32\times10^{-7} $m
    \item Wavelength of red light $\lambda_r=650$nm$=6.5\times10^{-7} $m
\end{itemize}
We know,
\begin{itemize}
    \item Planck constant $h=6.626\times10^{-34}$ J s
    \item $1$ eV $=1.602\times10^{-19}$ J
\end{itemize}
Energy of red light photon, $E=_r\frac{hc}{\lambda_r}=\frac{6.626\times10^{-34}\times3\times10^8}{6.5\times10^{-7}}$ J $=3.058\times10^{-19}$ J $=1.9$ eV \\
Energy of green light photon, $E=_r\frac{hc}{\lambda_r}=\frac{6.626\times10^{-34}\times3\times10^8}{5.32\times10^{-7}}$ J $=3.7365\times10^{-19}$ J $=2.3$ eV \\\\
\textcolor{cyan}{\textbf{The Energy of Red and Green light photon are 1.9 eV and 2.3 eV respectively.}}\\

\begin{center}
	\section*{\textcolor{goldColor}{\textbf{Problem C\\(a)}}}
\end{center}
Given,
\begin{itemize}
    \item Work function $W=2.3$ eV $=3.685\times10^{-19}$ J
\end{itemize}
We know,
\begin{itemize}
    \item Planck constant $h=6.626\times10^{-34}$ J s
\end{itemize}
If $f=$the minimum frequency of a photon to make an electron leave from a metal surface, therefore : $W=hf$\\
$\therefore f=\frac{W}{h}=\frac{3.685\times10^{-19}}{6.626\times10^{-34}}$ Hz $=5.56\times10^{14}$ Hz\\\\
\textcolor{cyan}{\textbf{The minimum frequency of a photon is $\mathbf{5.56\times10^{14}}$ Hz}}

\begin{center}
    \section*{\textcolor{goldColor}{\textbf{(b)}}}
\end{center}
Given,
\begin{itemize}
    \item Work function $W=2.3$ eV $=3.685\times10^{-19}$ J
    \item Wavelength of Photon $\lambda=2500$ \AA $=2.5\times10^{-7}$ m
\end{itemize}
We know,
\begin{itemize}
    \item Planck constant $h=6.626\times10^{-34}$ J s
    \item Speed of light $c=3\times10^8$ m/s
\end{itemize}
Energy of Photon, $E=\frac{hc}{\lambda}=\frac{6.626\times10^{-34}\times3\times10^8}{2.5\times10^{-7}}$ J $=7.95\times10^{-19}$ J $=4.96$ eV\\
The maximum kinetic energy $K_{max}=E-W=4.96-2.3$ eV $=2.66$ eV $\approx 2.7$ eV\\\\
\textcolor{cyan}{\textbf{The maximum kinetic energy of the emitted photoelectron is 2.7 eV}}

\newpage

\section*{\centering \textcolor{goldColor}{\textbf{Problem D\\ (a)}}}

Given,\\
\hspace*{0.5cm} A photoelectron moving at $v=5\times10^5$ m/s\\
We know,\\
\hspace*{0.5cm} Mass of electron $m=9.11\times10^{-31}$ kg\\\\
We estimate the uncertainty of velocity $\Delta v\approx v$  \\\\
Then, the uncertainty of momentum is, $\Delta p=mv$\\
According to  position–momentum uncertainty principle, $\Delta x \Delta p \geq \frac{h}{4\pi}$ \\\\
Therefore, $\Delta x \geq \frac{h}{4\pi\times mv}=\frac{h}{4\pi\times 9.11\times10^{-31}\mathrm{kg}\times5\times10^5\mathrm{m/s}}$\\
$\therefore \Delta x \geq 10^{-10}$ m\\\\
\textcolor{cyan}{\textbf{The estimated uncertainty of position is equal or greater than roughly }$\mathbf{10^{-10}}$\textbf{ m}}

\begin{center}
    \section*{\textcolor{goldColor}{\textbf{(b)}}}
\end{center}
Given,\\
\hspace*{0.5cm} A person is moving at $3$ m/s\\
\hspace*{0.5cm} Mass of person $m=75$ kg\\\\
We estimate the uncertainty of velocity $\Delta v\approx v$\\\\
Then, the uncertainty of momentum is, $\Delta p=mv$\\
According to  position–momentum uncertainty principle, $\Delta x \Delta p \geq \frac{h}{4\pi}$ \\\\
Therefore, $\Delta x \geq \frac{h}{4\pi\times mv}=\frac{h}{4\pi\times 75\mathrm{ kg}\times3\mathrm{  m/s}}$\\
$\therefore \Delta x \geq 2.3\times10^{-37}$ m\\\\
\textcolor{cyan}{\textbf{The estimated uncertainty of position is equal or greater than }$\mathbf{2.3\times10^{-37}}$\textbf{ m}}\\\\

\begin{center}
    \section*{\textcolor{goldColor}{\textbf{ Problem E }}}
\end{center} 

Quantum computer can be used to do very complex works. It can be used for complex calculations, simulations, AI, Data science, crypto, modeling etc. It works faster than normal computers. The two fields where scientists hope quantum computers will be especially useful, are :

\begin{enumerate}
    \item \textbf{Climate Modeling : }Quantum Computer can be used to Simulate atmospheric and oceanic changes to predict the weather. It can track and map precisely.
    \item \textbf{Nanotechnology : }Using Quantum Computer, we can simulate and design new Nanomaterials. It can also revolutionize Nanomedicine and Nano-electronics.
\end{enumerate}
\textcolor{cyan}{\textbf{Among many fields,Climate Modeling and Nanotechnology are two fields where scientists hope quantum computers will be especially useful.}}\\

Laptop and Quantum Computer have different mathematical system. They are not for same tasks. Quantum Computer is also very expensive. Furthermore, quantum hardware is completely unsuited for daily works because it cannot naturally process standard data formats like text or images, operates with high error rates, and requires massive cooling systems running near absolute zero.\\\\
\textcolor{cyan}{\textbf{So, today’s laptops are still better than quantum computers for most everyday tasks becouse of being economical, simple mechanism and suitable work daily works.}}


\end{document}
